\section{Related Work: Agent Safety and Attacks}
\label{sec:related}

\paragraph{Indirect prompt injection on the agentic Web.} The defining attack channel of tool-using agents is the observation itself: instructions hidden in tool results, web pages, or documents steer the agent while the user sees nothing~\cite{wang2026landscape}. The line of work is advancing quickly. Automatic generators such as PIA produce task-agnostic injection suffixes that transfer across scenarios~\cite{liu2024automatic}; WebInject exploits the visual channel of web agents~\cite{wang2025webinject}; WAInjectBench systematically benchmarks injection detections for web agents and finds detectors tuned on one page layout fail on others~\cite{liu2025wainjectbench}; and the newest work pushes toward covert payloads that keep the victim's visible behavior clean while exfiltrating data~\cite{lee2026will}. A parallel line attacks tool selection itself, causing the agent to invoke attacker-favored tools~\cite{shi2025prompt}. What these families share, and what motivates our benchmark, is that success is only fully visible in the resulting \emph{action stream}: which tools were called, with which arguments, feeding which downstream calls.

\paragraph{Malicious use and safety awareness.} A complementary threat model drops the adversary-in-the-middle and asks what happens when the user is the adversary. AgentHarm measures the harmfulness of agents that comply with malicious requests across 44 tool-enabled tasks and shows chat-level refusal training fails to transfer to agentic settings~\cite{andriushchenko2024agentharm}; ToolSword traces safety issues across planning, execution, and reflection stages of tool learning~\cite{ye2024toolsword}; R-Judge benchmarks whether models are even \emph{aware} of safety risk in recorded multi-turn agent interactions~\cite{yuan2024rjudge}; and Unsafer extends the audit to multi-turn tool-using agents, showing risk compounds across turns and proposing step-level defenses~\cite{li2026unsafer}. AgentSecBench aggregates prompt injection, privacy leakage, and tool-use integrity into one measurable surface~\cite{alpay2026agentsecbench}, and HarnessRisk evaluates the agent \emph{harness}---the layer managing tools, permissions, and state---over its lifecycle~\cite{bai2026harnessrisk}. These suites define \emph{what} to detect; our benchmark defines \emph{how} detection should be measured at the trajectory level, and imports their labeled traces as anomaly sources rather than re-inventing scenarios.

\paragraph{Defenses and their audit blind spot.} Defenses cluster into three designs. Input-layer hardening restructures queries or optimizes model preferences to resist injection~\cite{chen2024struq,chen2024secalign}. Plan- and runtime-layer defenses constrain \emph{what the agent may do}: IPIGuard checks tool calls against a tool-dependency graph derived from the plan~\cite{an2025ipiguard}, ToolSafe installs a proactive step-level guardrail with feedback~\cite{mou2026toolsafe}, AgentDoG provides a diagnostic guardrail framework for agent safety~\cite{liu2026agentdog}, NEXUS enforces structured runtime safety for tool use~\cite{hossain2026nexus}, AgentGuard repurposes the orchestrator for safety evaluation of tool orchestration~\cite{chen2025agentguard}, and WebMCP-Phalanx enforces trust boundaries for browser-integrated agents~\cite{lee2026webmcpphalanx}. All of these are step- or plan-scoped. Notably, even the most structural of them, IPIGuard, checks calls against the \emph{intended} plan rather than learning what anomalous execution \emph{looks like} across a corpus of real traces; none learns a detector over observed trajectory structure. Our benchmark is precisely the substrate such learned structural detectors require, and our disclosure ablation quantifies how much signal each defense granularity leaves on the table.

\paragraph{Benchmark validity and the closest work.} Two recent audits sharpen the stakes. Passing-the-Test shows prompt-injection detectors evaluated on the distributions they were tuned on lose most of their advantage out of distribution~\cite{liu2026passing}, and a validity audit of agent-safety benchmarks documents construct-validity failures---scores that reflect capability rather than safety~\cite{wang2026safety}---echoing concerns raised for agent evaluation generally~\cite{mohammadi2025evaluation}. On the data side, AutoDojo generates fresh prompt-injection scenarios to escape static test sets~\cite{ma2026autodojo}, and PASTABench assesses agents over \emph{sequential trajectories}~\cite{sun2026pastabench}. The closest work to ours is the AgentDrift/DriftNet pair: AgentDrift publishes step-labeled injection-hijacked trajectories, and DriftNet is a dual-head trajectory transformer that detects and localizes hijacks from the sequence of steps~\cite{pinjari2026agentdrift,pinjari2026driftnet}. Our contribution is orthogonal and composable with theirs: we compile trajectories into heterogeneous provenance graphs, generalize the anomaly definition beyond injection to malicious use and over-privilege, and evaluate \emph{graph-level} detectors under generalization, transfer, and mimicry controls that sequence-based evaluation does not exercise. To our knowledge, no prior benchmark treats the agent trajectory as a provenance graph or audits it with graph-level anomaly detection.
